(1) 作用在 \(u_n\) 的两根弹簧分别连向 \(v_n,v_{n-1}\)，而 \(v_n\) 连向 \(u_n,u_{n+1}\)，故
\[
m_1\ddot u_n=k(v_n+v_{n-1}-2u_n),\qquad
m_2\ddot v_n=k(u_n+u_{n+1}-2v_n).
\]
代入 \(u_n=Ue^{i(nqa-\omega t)}\)、\(v_n=Ve^{i(nqa-\omega t)}\)，得到系数矩阵
\[
\begin{pmatrix}2k-m_1\omega^2&-k(1+e^{-iqa})\\-k(1+e^{iqa})&2k-m_2\omega^2\end{pmatrix}\binom UV=0.
\]
行列式为零，即
\[
m_1m_2\omega^4-2k(m_1+m_2)\omega^2+4k^2\sin^2(qa/2)=0.
\]
两根正平方频率为
\[
\omega_\pm^2(q)=\frac{k}{m_1m_2}\left[m_1+m_2\pm\sqrt{(m_1+m_2)^2-4m_1m_2\sin^2(qa/2)}\right].
\]

(2) \(q=0\) 时 \(\omega_-(0)=0\)，是同相平移极限的声学支；\(\omega_+^2(0)=2k(m_1+m_2)/(m_1m_2)\)，为光学支。区边界 \(\sin^2(qa/2)=1\)，故
\[
\omega_-^2(\pi/a)=\frac{2k}{\max(m_1,m_2)},\qquad
\omega_+^2(\pi/a)=\frac{2k}{\min(m_1,m_2)}.
\]

(3) 因两支均随 \(\sin^2(qa/2)\) 单调变化，声学支上界是 \(\sqrt{2k/\max(m_1,m_2)}\)，光学支下界是 \(\sqrt{2k/\min(m_1,m_2)}\)。若质量不同，两端严格分离，中间无传播本征频率，即出现禁带。\(m_1=m_2\) 时两频率在折叠的区边界相遇，禁带闭合；这不是“无波动”，而是原胞可进一步缩小的结果。

(4) 在 \(q\to0\) 时 \(\sin(qa/2)\sim qa/2\)。将二次根式对 \(q\) 展开可得
\[
\omega_-^2(q)=\frac{k a^2}{2(m_1+m_2)}q^2+O(q^4),\qquad
c_s=\lim_{q\to0^+}\omega_-(q)/q=a\sqrt{\frac{k}{2(m_1+m_2)}}.
\]
若两质量相等为 \(m\)，则 \(c_s=(a/2)\sqrt{k/m}\)，与粒子间实际间隔 \(a/2\) 的单原子弹簧链声速完全一致。
